\section{Conclusion}
\label{sec:conclusion}

The six experiments in Track~E constitute the most systematic algorithmic
exploration of the congressional representation design space to date.  By
holding the core algorithm constant and varying design dimensions in controlled
ways, the experiments produce comparable, quantifiable evidence about the
costs and benefits of structural alternatives to single-member geographic
districting.

The five lessons we derive from this evidence are:

\begin{enumerate}
  \item \textbf{Geography is destiny for single-member districts.}
        The Rodden sorting gap is not a drawing artifact.  No single-member
        geographic drawing method --- algorithmic, commission-based, or
        judicially supervised --- eliminates systematic underrepresentation
        of Democrats relative to national vote share.  The cause is where
        Democratic voters live, not how district boundaries are drawn.

  \item \textbf{Multi-member districts partially but not fully escape sorting.}
        Three-member districts reduce the gap by 40\%; five-member by 65\%.
        The residual gap reflects extreme urban concentration that no
        realistic district size can dilute given state-boundary constraints.
        Multi-member districts also improve minority representation by 18\%
        without requiring geographic concentration.

  \item \textbf{Eliminating districts trades proportionality for local
        accountability.}  County representation improves proportionality in
        competitive states but introduces fractional voting, reduces
        constituent-representative bonds, and is constitutionally unavailable
        in its strict form without amendment.

  \item \textbf{Crossing state lines produces compact but fragmented districts.}
        National redistricting gains 12\% mean Polsby--Popper but loses 40\%
        of county preservation and creates 87 cross-state districts that
        fragment state-based political accountability.  The federalism cost
        is real but modest; the representation fragmentation cost is harder
        to quantify but significant.

  \item \textbf{The compactness--proportionality frontier is steep.}
        Gaining 1~pp of proportionality costs 0.015 PP units.  Achieving
        near-perfect proportionality requires maps no more compact than
        the commission-drawn maps that algorithmic redistricting is
        intended to improve.  No feasible single-election geographic system
        escapes this trade-off.
\end{enumerate}

These lessons support a clear design hierarchy.  For single-member systems
operating within current constitutional constraints, the optimal constraint
ordering is population equality, contiguity, and compactness --- with
community preservation as a soft secondary criterion.  Proportionality should
not be added as a drawing criterion within single-member systems because it
addresses a structural problem through a drawing instrument and produces
inferior results on both compactness and proportionality relative to structural
reform.

The structural reform that the evidence most clearly supports is multi-member
districts with proportional allocation via ranked-choice voting.  This reform
is constitutionally available without amendment, produces material improvement
in both proportionality and minority representation, and is the only mechanism
within a geographic representation system that partially escapes the Rodden
sorting constraint.  Its costs --- ballot complexity, weaker local ties ---
are real but normative rather than empirical, and the E-track evidence provides
the empirical foundation on which normative debates should proceed.

\subsection{Limitations and Future Research}

This paper draws on simulated algorithmic outcomes rather than actual electoral
results.  The vote shares used to compute proportionality deviations are
historical averages; future elections may differ.  The geographic data are from
the 2020 Census; demographic and residential sorting patterns will change over
time.  The sorting-gap estimates are derived from House election results and
may not generalize to other office levels.

Future research should examine hybrid systems --- multi-member districts in
urban areas combined with single-member districts in rural areas --- which are
not included in the E-track experiments.  Such hybrid systems might achieve
a better proportionality--compactness balance than either pure multi-member
or pure single-member systems, since they concentrate the multi-member
mechanism where it provides the most benefit (urban sorting concentrations)
while preserving the local-accountability advantages of single-member districts
where they are most valued (rural and small-town constituencies).

The international evidence (E.6) suggests that algorithmic redistricting
provides compactness improvements in all Westminster-system countries studied.
Future research should extend this analysis to mixed-member proportional
systems (Germany, New Zealand) and proportional representation systems (Nordic
countries) to test whether any of the lessons derived from the single-member
US context apply in fundamentally different electoral architectures.

\subsection{Coda: The Value of Algorithmic Counterfactuals}

The E-track research program demonstrates that algorithmic implementation of
constitutionally unavailable systems (national redistricting) and politically
controversial systems (partisan-similarity clustering) has significant
analytical value even when the systems cannot be adopted.  The ability to
implement and measure the outcomes of counterfactual systems is a uniquely
algorithmic contribution to the representation design literature.  Manual
redistricting cannot implement national redistricting (it is unconstitutional),
cannot implement strict partisan-fairness allocation (it is legally prohibited),
and cannot implement county representation at scale (it has never been attempted
in the US context).  Algorithmic redistricting can implement all of these as
measurement instruments, producing quantitative evidence about what we pay for
the constitutional and institutional choices we have made.

The lesson for political science methodology is that algorithmic implementation
of counterfactual institutions is a powerful tool for causal inference about
representation system design.  The lesson for redistricting reform is that the
evidence thus produced should inform the normative debates that determine which
institutional choices we make going forward.
